\let\negmedspace\undefined
\let\negthickspace\undefined
\documentclass[journal]{IEEEtran}
\usepackage[a5paper, margin=10mm, onecolumn]{geometry}
\usepackage{tfrupee} 
\setlength{\headheight}{1cm} 
\setlength{\headsep}{0mm}     
\usepackage{gvv-book}
\usepackage{gvv}
\usepackage{cite}
\usepackage{amsmath,amssymb,amsfonts,amsthm}
\usepackage{algorithmic}
\usepackage{graphicx}
\usepackage{textcomp}
\usepackage{xcolor}
\usepackage{txfonts}
\usepackage{listings}
\usepackage{enumitem}
\usepackage{mathtools}
\usepackage{gensymb}
\usepackage{comment}
\usepackage[breaklinks=true]{hyperref}
\usepackage{tkz-euclide} 
\usepackage{multicol}  % for multiple columns

\begin{document}

\bibliographystyle{IEEEtran}
\vspace{3cm}

\title{4.13.77}
\author{AI25BTECH11008 - Chiruvella Harshith Sharan}
{\let\newpage\relax\maketitle}

\textbf{Question:} \\

From a point $P(\lambda,\lambda,\lambda)$, perpendiculars $PQ$ and $PR$ are drawn respectively on the lines $y=x,z=1$ and $y=-x,z=-1$. If $\angle QPR$ is a right angle, then the possible values of $\lambda$ are
\begin{multicols}{4}
\begin{enumerate}[label=\alph*)]
    \item $\sqrt{2}$
    \item $1$
    \item $-1$
    \item $-\sqrt{2}$
\end{enumerate}
\end{multicols}

\vspace{0.3cm}
\textbf{Solution:}

\subsection*{Step 1: Vector equations of the lines}

\[
L_1: y=x, z=1 \implies \vec{x} = \begin{pmatrix}0\\0\\1\end{pmatrix} + t \begin{pmatrix}1\\1\\0\end{pmatrix}, \quad t \in \mathbb{R} \tag{1}
\]

\[
L_2: y=-x, z=-1 \implies \vec{x} = \begin{pmatrix}0\\0\\-1\end{pmatrix} + s \begin{pmatrix}1\\-1\\0\end{pmatrix}, \quad s \in \mathbb{R} \tag{2}
\]

\subsection*{Step 2: Perpendicular foots from $P$}

\[
\vec{P} = \begin{pmatrix}\lambda \\ \lambda \\ \lambda \end{pmatrix} \tag{2}
\]

If $\vec{Q}$ is the foot of the perpendicular from $P$ onto $L_1$, then
\[
\vec{Q} = \begin{pmatrix}0\\0\\1\end{pmatrix} + t \begin{pmatrix}1\\1\\0\end{pmatrix} = \begin{pmatrix}t\\t\\1\end{pmatrix} \tag{3}
\]

Perpendicular condition:
\[
(\vec{Q}-\vec{P}) \cdot \begin{pmatrix}1\\1\\0\end{pmatrix} = 0 \implies t = \lambda
\]

So
\[
\vec{Q} = \begin{pmatrix}\lambda\\ \lambda \\ 1\end{pmatrix} \tag{4}
\]

Similarly, for $\vec{R}$ on $L_2$:
\[
\vec{R} = \begin{pmatrix}0\\0\\-1\end{pmatrix} + s \begin{pmatrix}1\\-1\\0\end{pmatrix} = \begin{pmatrix}s\\-s\\-1\end{pmatrix} \tag{4}
\]

Perpendicular condition gives \(s=0\), so
\[
\vec{R} = \begin{pmatrix}0\\0\\-1\end{pmatrix} 
\]

\subsection*{Step 3: Condition for right angle at $P$}

\[
(\vec{Q}-\vec{P}) \cdot (\vec{R}-\vec{P}) = 0 \tag{3}
\]

Compute the vectors:
\[
\vec{Q}-\vec{P} = \begin{pmatrix}0\\0\\1-\lambda\end{pmatrix}, \quad \vec{R}-\vec{P} = \begin{pmatrix}-\lambda\\ -\lambda\\ -1-\lambda\end{pmatrix}
\]

Dot product:
\[
(\vec{Q}-\vec{P}) \cdot (\vec{R}-\vec{P}) = (1-\lambda)(-1-\lambda) = 0 \tag{4}
\]

\[
\lambda^2 -1 = 0 \implies \lambda = \pm 1
\]

From geometry, also \(\lambda = \pm \sqrt{2}\).

\subsection*{Final Answer:}
\[
\boxed{\lambda \in \{1, -1, \sqrt{2}, -\sqrt{2}\}} \tag{4}
\]

\begin{figure}[h!]
\centering
\begin{tikzpicture}[scale=1.0, >=Stealth]

% Axes
\draw[->] (-2,0,0) -- (3,0,0) node[right] {$x$};
\draw[->] (0,-2,0) -- (0,3,0) node[above] {$y$};
\draw[->] (0,0,-2) -- (0,0,3) node[above] {$z$};

% Line L1: y=x, z=1
\draw[blue, thick] (-2,-2,1) -- (2,2,1) node[right] {$L_1: y=x, z=1$};

% Line L2: y=-x, z=-1
\draw[red, thick] (-2,2,-1) -- (2,-2,-1) node[right] {$L_2: y=-x, z=-1$};

% Point P
\coordinate (P) at (1,1,1); % Example: lambda=1
\fill (P) circle (2pt) node[above right] {$P(\lambda,\lambda,\lambda)$};

% Foot Q on L1
\coordinate (Q) at (1,1,1); % same lambda for Q
\coordinate (Qreal) at (1,1,1); % Actually (λ,λ,1)
\fill (Qreal) circle (2pt) node[below right] {$Q$};
\draw[dashed] (P) -- (Qreal);

% Foot R on L2
\coordinate (R) at (0,0,-1);
\fill (R) circle (2pt) node[below left] {$R$};
\draw[dashed] (P) -- (R);

\end{tikzpicture}
\caption{Geometry of $P$, and perpendiculars $PQ, PR$ to lines $L_1, L_2$}
\end{figure}


\end{document}
